\section*{Chapter \ref{ch:b3:efficiency-cost-of-taxation} --- The Efficiency Cost of Taxation}

\begin{solution}{exo:b3:efficiency-cost-of-taxation:1}
In Britain the share went from 9.1~\% in 1913 to 19.9~\% in 1920 and 37.4~\% in 1946; in the
United States from 1.8~\% to 7.5~\% and 17.3~\%: steps at the two wars. Japan's largest
ten-year rise ended in 1944. Sweden's and Italy's came in peacetime: Italy's in the ten
years to 1983 (12.9 points); Sweden's, with the 1996 spike smoothed, in the ten years
to 2001 (11.1 points), across the jump of the mid-1990s, after a long rise from 15~\%
in 1950 to 29~\% in 1985. After 1950 only Japan fell below its 1913 share (to 10.1~\% at
the lowest, against 13.6~\%). % figdata: revenue, revenue_points
\end{solution}

\begin{solution}{exo:b3:efficiency-cost-of-taxation:2}
\omterm{def:b3:efficiency-cost-of-taxation:harberger}{Excess burden} $\tfrac12 \times 0.5 \times 0.2^2 \times 100 = 1$ million dalens. Revenue $0.2 \times 100 \times (1 - 0.5 \times 0.2) = 18$
million. Ratio $1/18 = 5.6~\%$.
\end{solution}

\begin{solution}{exo:b3:efficiency-cost-of-taxation:3}
\omterm{def:b3:efficiency-cost-of-taxation:mcpf}{Marginal excess burden} $0.1/(1 - 0.2) = 0.125$; \omterm{def:b3:efficiency-cost-of-taxation:mcpf}{marginal cost of public funds} 1.125. The
project returns 1.1, less than 1.125: not if it is financed by this tax.
\end{solution}

\begin{solution}{exo:b3:efficiency-cost-of-taxation:4}
Consumer prices that rise by less than the tax when the rate goes up, with
producer prices or margins falling; prices that respond differently to rises and
cuts of the same size; and smaller pass-through in markets where supply is less
elastic.
\end{solution}

\begin{solution}{exo:b3:efficiency-cost-of-taxation:5}
Set spending to 1 for each good, so revenue must be 0.2. The rule gives $0.2\tau_1 = 0.8\tau_2$, so
$\tau_1 = 4t$ and $\tau_2 = t$. Revenue: $4t(1 - 0.8t) + t(1 - 0.8t) = 5t - 4t^2 = 0.2$, so $t = (5 - \sqrt{21.8})/8 =
0.0414$: $\tau_1 = 16.5~\%$, $\tau_2 = 4.1~\%$. \omterm{def:b3:efficiency-cost-of-taxation:harberger}{Excess burden} $\tfrac12(0.2 \times 0.1655^2 + 0.8 \times 0.0414^2) = 0.00342$.
Uniform: $2\tau - \tau^2 = 0.2$ gives $\tau = 1 - \sqrt{0.8} = 10.6~\%$ and a burden $\tfrac12 \times 1.0 \times 0.1056^2 = 0.00557$. The
Ramsey rates cut the burden by about 39~\%.
\end{solution}

\begin{solution}{exo:b3:efficiency-cost-of-taxation:6}
$1/(1 + 1.5 \times 0.25) = 72.7~\%$; with $a = 2$, $1/(1 + 0.5) = 66.7~\%$. A larger $a$ means fewer very high
incomes above the threshold: the income on which a higher rate raises revenue,
$z_m - \bar z$, is small next to the income $z_m$ whose response it loses.
\end{solution}

\begin{solution}{exo:b3:efficiency-cost-of-taxation:7}
The net wage falls by $0.6/0.8 \times 10 = 7.5~\%$; the employer's cost rises by $0.2/0.8 \times 10 = 2.5~\%$.
Levied on workers, the outcome is the same to first order: the less elastic supply
side bears three quarters.
\end{solution}

\begin{solution}{exo:b3:efficiency-cost-of-taxation:8}
The tested solution is \code{ex08_composition.py}, in the chapter's \code{exercises/}
folder. Between 1990 and 2024, social contributions gained 25.1 points in Korea, the
largest change, while goods and services fell by 18.4 and trade by 10.7; France
moved 12.0 points toward income taxes; Germany's 1990 social contributions are
missing. % output: ex08_composition
\end{solution}

\begin{solution}{exo:b3:efficiency-cost-of-taxation:9}
Bunching in hours worked, or in earnings reported by employers rather than by the
taxpayers themselves; bunching among wage earners who cannot easily retime or
relabel income; and earnings that stay at the kink in later years instead of
returning to the smooth density.
\end{solution}

\begin{solution}{exo:b3:efficiency-cost-of-taxation:10}
The Lagrangian of the text, with $n$ goods, gives $\varepsilon_i\tau_i(1 + 2\lambda) = \lambda$ for each $i$, so $\tau_i\varepsilon_i$ is
the same across goods. With cross-price effects, the \omterm{def:b3:efficiency-cost-of-taxation:harberger}{excess burden} includes terms in
$\tau_i\tau_j$, and the first-order conditions mix the rates: a good whose consumption rises
when other goods are taxed should be taxed more than its own elasticity says. Leisure
cannot be taxed; a good bought with leisure (a complement to it) should be taxed
more, and a good that substitutes for leisure --- one bought by people who work
more, such as child care --- should be taxed less.
\end{solution}

\begin{solution}{exo:b3:efficiency-cost-of-taxation:11}
The marginal buncher's income response is the excess mass over the density, $120/0.05 =
2{,}400$ dalens. The elasticity is $(2{,}400/40{,}000)/\ln(0.80/0.65) = 0.06/0.208 = 0.29$.
\end{solution}

\begin{solution}{exo:b3:efficiency-cost-of-taxation:12}
With perfectly elastic supply the price paid rises by the full tax: buyers bear it
all. Demand falls by $1 \times 0.4 \times 0.1 = 4~\%$, and the lost units were worth to buyers between $p$
and the perceived price $1.04p$: the \omterm{def:b3:efficiency-cost-of-taxation:harberger}{excess burden} is $\tfrac12 \times 4 \times 0.04 = 0.08$, against $\tfrac12 \times 1 \times 0.1^2 \times 100 =
0.5$ if posted. The low salience lowers the triangle but leaves buyers spending as if
richer than they are, a cost the triangle does not show.
\end{solution}

\begin{solution}{pb:b3:efficiency-cost-of-taxation:brenland}
% figdata: brenland_stats
\textbf{1.} At 25{,}000: $0.20 \times 13{,}000 = 2{,}600$, an average rate of 10.4~\%. At 60{,}000: $0.20 \times 28{,}000 + 0.35 \times
20{,}000 = 12{,}600$, an average rate of 21.0~\%.

\textbf{2.} 20~\% and 35~\%.

\textbf{3.} 311 households (15.6~\%), with mean income 53{,}566 dalens.

\textbf{4.} 1{,}481{,}630 dalens, 21.8~\% of the 6{,}782{,}180 raised in all.

\textbf{5.} $a = 53{,}566/13{,}566 = 3.95$.

\textbf{6.} Four observations give no stable estimate; the rate on income from
40{,}000 to 100{,}000 is the one that most households in the upper bands face.

\textbf{7.} One point on $\sum(z_i - 40{,}000)$: 42{,}191 dalens.

\textbf{8.} $0.25 \times \sum\tau_iz_i/(1 - \tau_i) \times 0.01 = 22{,}740$ dalens.

\textbf{9.} $42{,}191/(42{,}191 - 22{,}740) = 2.17$.

\textbf{10.} $0.35 \times 0.25 \times 3.95/0.65 = 0.53$, so $1/(1 - 0.53) = 2.13$; close, the difference coming from the
four households at 45~\%.

\textbf{11.} $1/(1 + 3.95 \times 0.25) = 50.3~\%$.

\textbf{12.} The behavioural loss doubles to 45{,}480, more than the mechanical gain of
42{,}191: raising the rates lowers revenue, and the formula's value ($-12.8$) has no
meaning as a cost. With $e = 0.5$ the bands would be past their peak, which is then
$1/(1 + 1.97) = 33.6~\%$.

\textbf{13.} $2 \times 42{,}191 = 84{,}383$ dalens.

\textbf{14.} $2 \times 22{,}740 = 45{,}480$ dalens.

\textbf{15.} $84{,}383 - 45{,}480 = 38{,}903$ dalens.

\textbf{16.} $38{,}903/46{,}787{,}130 = 0.083$ points: from 18~\% to about 18.08~\%.

\textbf{17.} $38{,}903 \times (2.17 - 1.25) = 35{,}754$ dalens, to first order.

\textbf{18.} The 311 upper-band households gain; every household pays slightly more
value-added tax, the poorer ones a larger share of their income.

\textbf{19.} Its base, all consumption, is broad and hard to move: people can retime
or relabel income far more easily than they can stop consuming.

\textbf{20.} No: 1.5 is less than 2.17.

\textbf{21.} Yes: 1.5 exceeds 1.25.

\textbf{22.} The behavioural loss is $e \times 90{,}960$ per point (22{,}740 at $e = 0.25$). The cost is 1.5 when
$e \times 90{,}960 = 42{,}191/3$: $e = 0.155$.

\textbf{23.} Reported incomes around past changes of the upper rates, against a group
whose rates did not change; bunching at the 40{,}000 and 100{,}000 kinks; and how much
of the response is retiming or shifting.

\textbf{24.} They range from 0.1--0.2 near the threshold to 0.83 in one year for the
whole group: the elasticity depends on who is taxed and how easily income can be
moved, and an announced rise invites retiming.

\textbf{25.} The elasticity measures how much income escapes the base; a broader base,
with fewer ways to shift income, lowers it and with it the cost of funds.

\textbf{26.} With an \omterm{def:b3:efficiency-cost-of-taxation:eti}{elasticity of taxable income} of 0.25, the \omterm{def:b3:efficiency-cost-of-taxation:mcpf}{marginal cost of public
funds} of Brenland's upper bands is 2.17 dalens per dalen raised.
\end{solution}
